\scanpage{12}
% archive sheet number: 11
 i) \uline{Formule d\deleted{'}e}\deleted{adjonction}\inserted{\uline{transposition ou d'adjonction}}, disant que Gysin \deleted{est} sur la cohomologie à supports propres est transposé de l'homomorphisme image inverse. J'ai un sentiment que cela résulte tautologiquement de la définition, \deleted{\unsure{et n'a} guère de sens que dans}

 j) \uline{Formules de \unsure{finéarité|linéarité}} (notamment pour $f$ une immersion de schémas relatifs lisses)\handwritten{,} nécessaire pour Jouanolou; regarder aussi la formule de linéarité pour variété\deleted{s} éclatée\deleted{s} déduite\deleted{s} de cette situation. Si possible, donner formulation commune ?

Comme plan possible du par.4, je vois

4.1. Définition et notations.

4.2. Formulaire général.

4.3. Formulaire spécial\deleted{\unsure{s|e}}, pour classes de cohomologie.

Du point de vue notations, je suggère, pour ne pas trainer indéfiniment des $d-d'$ dans les notations, de travailler systématiquement avec la \uline{dimension virtuelle} $d(f)=d-d'$ de $f$; c'est elle qui mérite une lettre simple, telle que $d$.

Dans 4.3, inclure le cas $f$ propre, $d=d'$; dire que dans ce \deleted{5. \uline{Classes fondamentales locales}} cas on a
\[
\mathrm{Tr}_f(1_X)=N\,1_Y
\]
\margin{\handwritten{\begin{tikzpicture}\draw[gray] (0,0) .. controls (-0.7,1) and (-0.7,2.5) .. (0,3.4);\end{tikzpicture}}}
où $N$ est une fonction entière localement constante sur $Y$, appelée \uline{degré} de $f$, et qu'on peut caractériser comme étant égale, sur \deleted{tout} le plus grand ouvert de $Y$ sur lequel $f$ soit fini localement libre, comme égale au degré habituel. Conclure que $\mathrm{Ker}(f^{\handwritten{*}}:\mathrm{H}^{\handwritten{\unsure{*|+}}}(Y,G)\handwritten{\longrightarrow}\mathrm{H}^{\handwritten{*}}(X,G_X))$ est annulé par $N$, donc\inserted{\handwritten{$f^*$}} est injectif (et même inversible à gauche) si $N=1$. \handwritten{[}Si $S=\mathrm{Spec}\,k$ et si $Y$ est irréductible, de point générique $\handwritten{\eta}$,\inserted{\handwritten{($d$ et $d'$ quelconques)}} et si $N$ est le pgcd des degrés des $0$-cycles de $X_{\handwritten{\eta}}$\inserted{\handwritten{$\neq\varnothing$}}, on peut de même trouver un homomorphisme\handwritten{,} fonctoriel en $G$\handwritten{\unsure{$\leftarrow$}}\handwritten{,} $\mathrm{H}^{\handwritten{*}}(X,G_X)\handwritten{\longrightarrow}\mathrm{H}^{\handwritten{*}}(Y,G)$ dont le composé avec l'homomorphisme image inverse est $N.\mathrm{id}$.\handwritten{]}
